\section{Experimental Evaluation}\label{sec:results}
\subsection{Objectives and Setup}
The experiments test whether the specified mechanics produce their expected
results. All six use synthetic inputs and fixed decision policies where
actions are required. They do not fit participant preferences or measure
real-market forecasting performance. The calculations and policies are small
enough to reconstruct from the descriptions below, and no model inference is
used in these six experiments.

Verification concerns whether the program executes the intended rules;
behavioral validation concerns whether those rules and decisions represent
real participants. This distinction determines the interpretation of every
result. We evaluate exact arithmetic, round timing, resource accounting,
replay, and a controlled removal of interaction mechanisms. We report a
separate small live API check only as integration evidence.

\subsection{Experiment 1: Incentives in a Coordination Game}
A hypothetical founder chooses whether to fund a review. A liquidity partner
chooses whether to commit resources to disclosed escrow or wait. Escrow means
resources are held under stated release conditions; it is part of the
fictional game's interpretation, not a deployed escrow service.
Table~\ref{tab:game} supplies both participants' utilities. Every number is an
assumed preference value, not cash or an estimated market payoff.

\begin{table}[t]
\caption{Synthetic utilities: (founder, partner).}
\label{tab:game}
\centering\footnotesize
\begin{tabular}{@{}lcc@{}}
\toprule & Partner commits & Partner waits \\
\midrule
Founder funds review & $(6,5)$ & $(-2,1)$ \\
Founder skips review & $(2,-4)$ & $(0,0)$ \\
\bottomrule
\end{tabular}
\end{table}

Enumeration finds two pure Nash profiles: review/commit and skip/wait.
In the first, either unilateral change lowers the deviator's utility; the
same is true in the second. The example therefore captures coordination
rather than a universally preferred individual action.

Given an assumed partner commitment probability of 0.4, the founder's
expected utility is $0.4\times6+0.6\times(-2)=1.2$ for review and
$0.4\times2+0.6\times0=0.8$ for skipping. These weighted averages are needed
to compare choices under the supplied uncertainty. Ignoring the probabilities
would compare isolated payoffs rather than the stated decision problem.
Skipping has regret 0.4 utility points. The result verifies the calculation,
not that 40\% of real partners would commit.

\subsection{Experiment 2: Updating Explicit Beliefs}
The two hypotheses are that attention persists or fades, with prior
probabilities 0.4 and 0.6. A first signal has likelihood 0.75 under persistence
and 0.25 under fading. Equation~\eqref{eq:bayes} gives a persistence posterior
of $2/3$. A second signal has corresponding likelihoods 0.6 and 0.3; assuming
conditional independence, the posterior becomes $4/5$.

The second update multiplies each hypothesis's prior by both signal
likelihoods and normalizes over the two hypotheses. Conditional independence
is what permits that multiplication. Without it, a joint likelihood would
be needed to avoid double-counting correlated signals. The resulting 0.8
checks sequential updating and normalization under supplied assumptions;
it is not a probability estimated from current news.

\subsection{Experiment 3: Timing, Accounting, and Replay}
Two builders each start with 100 cents. A hypothetical readiness metric
starts at one point and is bounded between zero and ten. Before round-one
observations, a scheduled event adds one point. Both builders then take a
60-cent action contributing six points each. Aggregate settlement reaches
14 and enforces the upper bound of ten. A round-two event subtracts one
before the next observation. The remaining 40 cents make building again
unaffordable, so both participants wait.

\begin{table}[t]
\caption{State timing in the synthetic two-builder world. Cash is in cents.}
\label{tab:timing}
\centering\footnotesize
\begin{tabular}{@{}lrrrr@{}}
\toprule Stage & Readiness & Alice & Bob & Spent \\
\midrule
Initialization & 1 & 100 & 100 & 0 \\
Round-one observation & 2 & 100 & 100 & 0 \\
Round-one settlement & 10 & 40 & 40 & 120 \\
Round-two observation & 9 & 40 & 40 & 120 \\
Terminal state & 9 & 40 & 40 & 120 \\
\bottomrule
\end{tabular}
\end{table}

Table~\ref{tab:timing} distinguishes settlement from the next round's opening
state. Build publications become visible to subscribers in round two.
Remaining cash plus spending is $40+40+120=200$, satisfying the accounting
proposition. With assumed cash weight one and readiness weight ten, each
participant's terminal utility is $40+10\times9=130$, compared with
$100+10\times1=110$ at declared initialization. This change includes the
scheduled events as well as actions. Replay reproduces the complete report
exactly. The build-then-wait policy is imposed, so the experiment tests
execution rather than Jev's ability to select useful actions.

\subsection{Experiment 4: Scoring Probabilistic Forecasts}
A forecast needs comparison with a resolved outcome; a high confidence field
alone cannot establish accuracy. For binary cases, let $p_j$ be the forecast
probability of a positive outcome, $y_j$ the resolved outcome (zero or one),
and $n$ the number of cases. These symbols are used only for forecast
evaluation, distinct from the participant action scores $q_{i,t}(a)$.
The implemented scores are
\begin{align}
\operatorname{Brier}&=\frac{1}{n}\sum_{j=1}^{n}(p_j-y_j)^2,
\label{eq:brier}\\
\operatorname{LogLoss}&=-\frac{1}{n}\sum_{j=1}^{n}
\big[y_j\log p_j+(1-y_j)\log(1-p_j)\big].
\label{eq:logloss}
\end{align}
Both scores are minimized by accurate probabilities in expectation
\cite{brier1950,gneiting2007}. Equation~\eqref{eq:brier} measures the average
squared gap between probability and outcome. Equation~\eqref{eq:logloss}
assigns a much larger penalty to confidently wrong forecasts. Division by
$n$ makes the scores averages rather than totals that grow with case count.
Without outcome-based scoring, one could rank forecasts by their confidence
or narrative appeal without testing whether they were right.

Two synthetic training outcomes imply a baseline positive rate of 0.5.
Two later synthetic cases have outcomes $(0,1)$ and supplied forecasts
$(0.25,0.75)$. Brier score is 0.0625 and log loss is 0.287682, compared with
0.25 and 0.693147 for the constant 0.5 baseline. Log loss bounds its numeric
inputs between $10^{-15}$ and $1-10^{-15}$ to avoid taking the logarithm of
zero. This numerical safeguard is not evidence of calibration. Two invented
cases verify the formulas, not forecasting skill. Chronology checks prevent
declared training records from following the evaluation cases, but cannot
authenticate the underlying data timestamps.

\subsection{Experiment 5: Communication and Joint Effects}\label{sec:ablation}
An ablation removes a mechanism while keeping the remaining experiment fixed.
Here, two binary interventions separate information delivery from a joint-action
effect. Alice and Bob each start with 100 cents, and readiness starts at zero.
Alice builds in round one, spends 20 cents, adds one readiness point, and
publishes a message. In round two Alice waits. Bob joins only after receiving
Alice's message, spending 30 cents and adding two points; otherwise he waits.
A declared joint rule adds three points when Alice waits and Bob joins.
Both participants follow these programmed policies in every condition.

\begin{table}[t]
\caption{Factorial ablation. Readiness is a hypothetical metric; cost and
remaining cash are cents. Initial total cash is 200 in every condition.}
\label{tab:ablation}
\centering\footnotesize
\begin{tabular}{@{}llrrr@{}}
\toprule Delivery & Joint rule & Readiness & Cost & Cash left \\
\midrule
On & On & 6 & 50 & 150 \\
On & Off & 3 & 50 & 150 \\
Off & On & 1 & 20 & 180 \\
Off & Off & 1 & 20 & 180 \\
\bottomrule
\end{tabular}
\end{table}

With delivery enabled, Bob joins. Removing only the joint rule leaves that
choice unchanged and reduces readiness from six to three. Removing delivery
makes Bob wait, so the joint rule cannot activate and readiness remains one.
Thus the mechanism that changes the selected action is distinguishable from
the mechanism that changes its numerical consequence. All four trajectories
replay exactly and reconcile their resources. These results describe effects
inside one stipulated program; they do not identify how real people respond
to information or establish that communication improves welfare.

\subsection{Experiment 6: Preferences Versus Persona Prose}
The eight-participant pilot in Section~\ref{sec:cast} tests whether declared
numeric preferences have executable consequences. Every condition uses a
local deterministic policy selecting the greatest immediate own-action utility
change, with stable menu-order tie breaking. This policy deliberately ignores
narrative motivations, opponent responses, and future effects. It establishes
a baseline whose limits are known, rather than attributing strategic reasoning
to an untested model.

Under the declared preferences, the participation-oriented participant joins
and the risk-sensitive participant waits. Changing only the former's risk
weight from minus 500 to minus 5,000 reverses its isolated join value from
8,500 to minus 5,000 utility points, and it waits instead. Across two rounds,
terminal participation changes from 44 to 34, risk from 24 to 18, and total
spending from 120,000 to 110,000 cents. Readiness remains 44 and credibility
98 in both conditions. Remaining cash plus spending equals 1,000,000 cents.
All trajectories replay exactly.

A third condition changes every motivation and goal to emphasize caution,
while holding numeric preferences and rules fixed. The local policy's actions
and numerical results remain identical to the original condition. Thus prose
alone does not implement a preference in this baseline, while changing the
implemented incentive changes both choices and outcomes. These are synthetic
mechanism checks, not results for Jev or observations of real participants.
Testing whether Jev follows, deviates from, or strategically improves on this
baseline requires separately recorded model experiments.

\subsection{Scale and Live Integration Checks}
Synthetic execution checks cover 10, 20, and 100 participants. An orchestration
check executes 100 participants over two rounds with 200 injected typed
responses, verifying message timing, resource constraints, and replay. Injected
responses are controlled test data, so this result does not measure live
100-participant latency, token usage, or behavioral quality.

A separate authenticated TypeSafe check requested one Choice decision using
synthetic state. A cautious participant was instructed to preserve capital
until verification was complete; verification was incomplete and the options
were wait or commit. The returned model was Jev 1.14.0 and the selected action
was wait, with scores one for wait and zero for commit. The response reported
409 input tokens and 31 output tokens; client-observed inference time was
0.314 seconds. At the cited list price the input estimate is \$0.000017178.
This single, deliberately simple request verifies connectivity and typed
response handling, not a full live simulation or calibrated certainty.

Agent interoperability was checked with the official Python MCP SDK 1.30.0
against the actual server subprocess, negotiating protocol version 2025-11-25
\cite{mcp2025}. The client discovered eleven tools, three resources, and one
construction prompt. It supplied a modified world, validated and executed it
with the offline policy, retrieved results, verified replay, and exported
Markdown. An installed-package check repeated the journey outside the source
checkout with zero provider calls. These tests establish a working protocol
path, rather than the quality of an agent's authored world or support for
every host application's configuration. The newer 2026 protocol is not claimed.

An operator-signed NEAR testnet registration has also been completed.
Its role is to exercise artifact registration, independently of model inference.
It does not establish external-user wallet acceptance, decentralized inference,
or the validity of a forecast. The optional NEAR AI Cloud author route has
not completed a funded generation. These operational checks are separate
from the six synthetic experiments.
